\subsection{On-premise řešení}
\label{subsec:deployment-on-premise}

On-premise nasazení znamená provozování jazykového modelu přímo na lokálních serverech, bez využití externích cloudových služeb.
V tomto případě je model hostován na serverové infrastruktuře a veškeré zpracování dat probíhá lokálně.

Výhodou on-premise řešení je především nezávislost na externích poskytovatelích služeb a absence poplatků za jednotlivé API požadavky.
Správci mají zároveň možnost plně kontrolovat konfiguraci modelu, způsob jeho aktualizace i způsob škálování výpočetních prostředků.
Na druhou stranu je nutné zajistit dostatečné hardwarové prostředky, zejména výkonné grafické akcelerátory (GPU), dostatek operační paměti a vhodnou infrastrukturu pro provoz a správu modelů.

Pro usnadnění provozu jazykových modelů v lokálním prostředí existuje několik softwarových nástrojů, které poskytují rozhraní pro jejich správu, spouštění a inferenci:

\begin{itemize}
    \item \textbf{Ollama}~\cite{ollama} -- nástroj zaměřený na jednoduché spouštění a správu jazykových modelů v lokálním prostředí.
    Poskytuje jednotné rozhraní pro práci s různými modely a umožňuje jejich snadné stažení, spuštění i aktualizaci.
    Součástí nástroje je také API, které umožňuje integraci modelů do externích aplikací.

    \item \textbf{Groq}~\cite{groq} -- platforma zaměřená na vysoce výkonnou inferenci jazykových modelů.
    Nabízí optimalizovanou infrastrukturu pro rychlé zpracování požadavků nad velkými modely.
    Přestože je Groq často využíván jako cloudová služba, jeho technologie je založena na specializovaném hardwaru navrženém pro efektivní provoz modelů.

    \item \textbf{vLLM}~\cite{vllm} -- knihovna určená pro efektivní provoz velkých jazykových modelů, která se zaměřuje zejména na optimalizaci paměťového využití a paralelního zpracování požadavků.
    Využívá techniky jako \textit{PagedAttention}, které umožňují efektivně sdílet paměť mezi více požadavky a tím výrazně zvyšují propustnost systému.

    \item \textbf{llama.cpp}~\cite{llamacpp} -- open-source implementace zaměřená na provoz modelů rodiny LLaMA a jejich variant i na běžných procesorech bez nutnosti využití GPU\@.
    Nástroj umožňuje provoz kvantizovaných verzí modelů, což výrazně snižuje paměťové nároky a umožňuje nasazení i na méně výkonném hardwaru.
\end{itemize}

Volba konkrétního nástroje závisí především na dostupném hardwaru, požadované rychlosti inferenčního procesu a způsobu integrace do existující infrastruktury systému Kelvin.
Zatímco nástroje jako \textit{vLLM} jsou vhodné pro provoz ve výkonném serverovém prostředí s GPU akcelerací, řešení jako \textit{llama.cpp} umožňují provoz modelů i na méně výkonném hardwaru.
Nástroj \textit{Ollama} pak poskytuje kompromis mezi jednoduchostí nasazení a možností integrace do aplikací prostřednictvím API\@.

\subsubsection*{Hardwarové nároky}

Pro on-premise nasazení je zcela zásadní dostupnost vhodného hardware.
Inference jazykových modelů je výpočetně náročná operace (viz \secref{subsec:llm-computational-cost}), proto je využití GPU s dostatečnou kapacitou paměti VRAM téměř nezbytné pro přijatelné výkony.
přibližné nároky na hardware pro různé modely lze vidět v \tabref{tab:memory-requirements}.

\TODO{Doplnit porovnání CPU vs GPU pro inferenci. Spustit ollama na argexu a změřit výkon pro různé modely.}

\subsubsection*{Limity}

On-premise nasazení přináší vlastní sadu omezení:

\begin{itemize}
    \item \textbf{Omezení dostupným hardware} -- kvalita a rychlost analýzy jsou přímou funkcí dostupného GPU.
    Na nedostatečném hardware je nutné volit menší nebo silněji kvantizované modely, což může snižovat kvalitu výstupů.

    \item \textbf{Škálovatelnost} -- cloudové API automaticky škáluje s nárůstem zátěže, kdežto lokální server má fixní kapacitu.
    Při hromadném odevzdávání se požadavky řadí do fronty a čekají na zpracování.

    \item \textbf{Správa modelu} -- nové verze modelů musí být ručně staženy a nasazeny.
    Aktualizace nebo výměna modelu vyžaduje přímou intervenci správce systému.
\end{itemize}

\subsubsection*{Ochrana dat}

Z pohledu ochrany dat je on-premise nasazení výrazně jednodušší situace.
Studentské zdrojové kódy a zadání úloh nikdy neopustí infrastrukturu školy.
Neexistuje třetí strana, které by musela být svěřena správa dat, a není potřeba uzavírat DPA s externím poskytovatelem.
Z hlediska GDPR jde o přímočarou variantu: správce dat (škola) je totožný s provozovatelem infrastruktury.

\subsubsection*{Náklady}

Náklady on-premise nasazení mají odlišnou strukturu než u cloudového API:

\begin{itemize}
    \item \textbf{Jednorázové pořizovací náklady} -- GPU server vhodný pro inferenci (např.\ NVIDIA RTX 4090 s 24 GB VRAM) se pohybuje v ceně od přibližně 60~000 Kč výše.
    Pro výkonnější varianty (NVIDIA A100) jsou náklady řádově vyšší (200~000--500~000 Kč).

    \item \textbf{Průběžné provozní náklady} -- elektřina (moderní GPU má spotřebu 150--400 W při zatížení), chlazení serverovny a lidská práce nutná pro správu a aktualizaci systému.

    \item \textbf{Nulové náklady za token} -- po zaplacení hardware jsou marginalní náklady na každé zpracované odevzdání prakticky nulové, bez ohledu na počet odevzdání.
\end{itemize}

Pro systém s malým počtem uživatelů a odevzdání může být on-premise pořizovací investice ekonomicky nevýhodná ve srovnání s cloudovým API\@.
Naopak pro systém se stovkami studentů zpracovávajícími tisíce odevzdání ročně se on-premise investice může vyplatit v horizontu 1--2 let.

\begin{table}[H]
    \centering
    \resizebox{\textwidth}{!}{%
        \begin{tabular}{lll}
            \toprule
            \textbf{Vlastnost}   & \textbf{Cloudové API}              & \textbf{On-premise nasazení} \\
            \midrule
            Hardwarové nároky    & Žádné (výpočet u poskytovatele)    & GPU server (10+ GB VRAM)     \\
            Pořizovací náklady   & Nulové                             & Vysoké (GPU hardware)        \\
            Průběžné náklady     & Za token (opakující se)            & Elektřina a správa           \\
            Škálovatelnost       & Automatická (rate limity)          & Fixní kapacita               \\
            Ochrana dat          & Data opouštějí školu               & Data zůstávají na škole      \\
            Kvalita modelů       & Nejnovější (GPT-4o, Claude Sonnet) & Závisí na hardware           \\
            Latence              & Závislá na internetu a zátěži API  & Závisí na hardware           \\
            Správa a aktualizace & Automatická (poskytovatel)         & Manuální (správce)           \\
            \bottomrule
        \end{tabular}
    }
    \caption{Srovnání cloudového API a on-premise nasazení jazykového modelu z hlediska klíčových vlastností}
    \label{tab:deployment-comparison}
\end{table}

\endinput